% Standalone section for inclusion in a report using amsmath, amssymb,
% amsthm and theorem/lemma/proposition environments.
% New guarded-core argument, developed against ConciseIPRA b17dd5487f34.

\section{A guarded core with a deterministic introduction rule}
\label{sec:guarded-core}

The change is to the public formula language and the introduction rule,
not to the numerical domain.  An atomic matrix immediately below a
paired binder must have a fixed, syntactically recognizable zero guard.
Such a matrix has a provable zero instance.  Consequently it needs only
the progress introduction; the separate atomic case introduction can
be deleted.  Compound matrices retain the case introduction.

\subsection{Complete syntax, axioms, and inference rules}

Numerical terms are those of the inherited scalar calculus:
\[
 t ::= x\mid 0\mid S(t)\mid I(n,a;u.h).
\]
Only the displayed variable $u$ is bound in $h$.  There is one numerical
sort, fixed finite term syntax, and capture-free substitution.  Put
\[
 C(v,a,b):=I(v,a;u.b),\qquad 1:=S0,
\]
where $u$ is fresh.  For a literal equation $E\equiv(s=t)$, define
\[
 G_v E\;:=\;\bigl(C(v,0,s)=C(v,0,t)\bigr).
\tag{Guard}\label{eq:guarded-guard}
\]
The displayed abbreviations expand into terms of the given syntax.

Call a literal equation a \emph{zero-initialized matrix in $y$} when it
has the syntactic form
\[
 B(x,y)\equiv
 I(y,0;u.h)=I(y,0;u.k).
 \tag{Canonical atomic matrix}\label{eq:guarded-canonical-matrix}
\]
Here $h,k$ may depend arbitrarily on the state variable $u$, on $x,y$,
and on fixed parameters; state binders are renamed as necessary.
The selector guard $G_yE$ is a special case, with state-independent
updates.  Every matrix of \eqref{eq:guarded-canonical-matrix} has the
theorem $B(x,0)$ by the two zero computations and equality.

The canonical formulas of $\mathcal G$ are generated as follows:
\begin{enumerate}
\item Every literal equation $s=t$ is a formula.
\item $\mathsf D_{x,y}B$ is a formula when $B$ is already a compound
      canonical formula (its root is $\mathsf D$).
\item $\mathsf D_{x,y}B(x,y)$ is a formula whenever $B$ has the
      zero-initialized shape \eqref{eq:guarded-canonical-matrix}.
\end{enumerate}
The bound variables $x,y$ are distinct.  Rule 3 is the restriction on
atomic matrices.  It does not require an equality decision procedure or
a proof that an arbitrary base is true: the guard has the displayed
literal shape.  The designated contradiction is $0=1$.

Retained contexts contain reusable clauses
$\langle A;y.B\rangle$, where $A,B$ are canonical formulas and
$y\notin\operatorname{FV}(A)$.  An ordinary assumption $F$ abbreviates
the constant clause $\langle H;v.F\rangle$, where $v$ is fresh and
$H$ is the fixed computation theorem $I(0,0;u.0)=0$.

The two nonlogical axiom schemas, available in every context, are
\[
 I(0,a;u.h)=a,\qquad
 I(Sn,a;u.h)=h[I(n,a;u.h)/u].
\tag{I0,IS}\label{eq:guarded-computation}
\]
Besides these, there are precisely the following rules:
\[
 \frac{\Gamma\vdash A}{\Gamma\vdash B(t)}
 \quad\text{if }\langle A;y.B(y)\rangle\in\Gamma
 \qquad(\mathrm{Use}),
\]
\[
 \frac{\Gamma\vdash s=t\qquad\Gamma\vdash F(s)}
      {\Gamma\vdash F(t)}\qquad(\mathrm{Tr}),
\]
\[
 \frac{\Gamma\vdash\mathsf D_{x,y}B(x,y)\qquad
       \Gamma,\langle B(a,0);y.B(a,y)\rangle\vdash F}
      {\Gamma\vdash F}\qquad(\mathrm{DE}),
\]
and the single, root-determined introduction schema
\[
 \frac{\Gamma,B(t,\eta_B(z))\vdash B(t,Sz)}
      {\Gamma\vdash\mathsf D_{x,y}B(x,y)},
 \qquad
 \eta_B(z)=
 \begin{cases}
 z,& B\text{ is a literal equation},\\
 0,& B\text{ is compound}.
 \end{cases}
 \tag{DI$_{\rm root}$}\label{eq:guarded-di}
\]
There is no certificate choice between these two cases.  At an atomic
matrix the guard condition on the conclusion is mandatory in the
canonical calculus.  In DE, $a$ is fresh for $\Gamma,F$ and the fixed
parameters.  In DI, $z$ is fresh for $\Gamma,t$ and the fixed
parameters.  The displayed local resources are discharged.
Transport ranges over canonical formula templates, and capture-free
numerical substitution preserves their guard shape.

There are no additional primitive logical rules, reflexivity axioms,
successor axioms, coverage axioms, or induction schemas.
Weakening, exchange, contraction, renaming, numerical substitution, and
ordinary-assumption cut are admissible by induction on derivations.
For cut, replace a use of the corresponding constant clause by the
supplied proof, renaming eigenparameters first.

\subsection{Padding, equation implication, and numerical cases}

Write $A\dashv\vdash B$ for interderivability under arbitrary additional
retained context.  The following preliminary arguments use no
contextual induction rule.

\begin{lemma}[Equality and guard endpoints]
Reflexivity, symmetry, transitivity, numerical congruence, and their
usual contextual uses are derivable.  Moreover,
\[
 \vdash G_0 E,\qquad G_{Sz}E\dashv\vdash E.
\]
\end{lemma}
\begin{proof}
For reflexivity, put $r=I(0,t;u.0)$.  The computation equation $r=t$
serves as both premises of transport with motive $v=t$, yielding
$t=t$.  Standard equality motives then give the remaining equality
rules.  The two guard claims follow from the two computation equations
for $C$, together with equality transport.
\end{proof}

For an equation $E$ let
\[
 L(E):=\mathsf D_{u,v}G_v E,
 \qquad
 P(A):=
 \begin{cases}L(A),&A\text{ is an equation},\\A,&A\text{ is compound}.
 \end{cases}
\tag{Padding}\label{eq:guarded-padding}
\]
The newly bound $u,v$ are fresh for $E$.  In particular, $P(A)$ is
always compound.

\begin{lemma}[Padding]
$L(E)\dashv\vdash E$, and therefore $P(A)\dashv\vdash A$.
\end{lemma}
\begin{proof}
Under $E$, atomic progress introduction derives $L(E)$: its target
$G_{Sz}E$ follows from the retained $E$, independently of the local
induction assumption.  Conversely, DE on $L(E)$ opens
$\langle G_0E;v.G_vE\rangle$.  Supply the guard-zero theorem and use
the clause at $1$; guard removal gives $E$.
\end{proof}

We will need to reason by numerical cases under an assumption that
contains the case variable.  It is essential not to violate the
eigenvariable condition by directly retaining that assumption during
case introduction.  A small implication macro resolves this issue.

For equations $U\equiv(s=t)$ and $V\equiv(r=q)$ put
\[
 M_w(U,V):=(C(w,s,r)=C(w,t,q)),\qquad
 U\Rightarrow_g V:=\mathsf D_{a,w}L(M_w(U,V)),
\tag{Equation implication}\label{eq:guarded-implication}
\]
with fresh padding variables and unused $a$.
The matrix of the outer binder is compound.

\begin{lemma}[Implication between equations]
The macros in \eqref{eq:guarded-implication} admit both rules
\[
 \frac{\Gamma,U\vdash V}{\Gamma\vdash U\Rightarrow_g V},
 \qquad
 \frac{\Gamma\vdash U\Rightarrow_g V\qquad\Gamma\vdash U}
      {\Gamma\vdash V}.
\]
\end{lemma}
\begin{proof}
The endpoint identities for $C$ and padding give
$L(M_0(U,V))\dashv\vdash U$ and
$L(M_{Sz}(U,V))\dashv\vdash V$.
Compound case DI therefore gives introduction.  DE, the first endpoint,
Use at $1$, and the second endpoint give elimination.  These endpoint
conversions use numerical transport with parameters outside the fresh
padding binders; they do not assume a binder-congruence theorem.
\end{proof}

\begin{lemma}[Numerical cases]
For every canonical $A(y)$ the following rule is admissible:
\[
 \frac{\Gamma\vdash A(0)\qquad\Gamma\vdash A(Sz)}
      {\Gamma\vdash A(t)},
 \qquad z\text{ fresh for }\Gamma\text{ and the fixed parameters}.
 \tag{NC}\label{eq:guarded-cases}
\]
\end{lemma}
\begin{proof}
First suppose $A$ is compound.  Case DI with an unused witness coordinate
introduces $\mathsf D_{u,y}A(y)$, ignoring its local base assumption in
the successor premise.  DE opens $\langle A(0);y.A(y)\rangle$; supply
the given base proof and use it at $t$.

If $A$ is an equation, apply this already proved compound case rule to
$L(A(y))$.  Padding converts the two premises to its endpoints and
converts the conclusion back.  No progress induction is being used in
this construction beyond the special padding introductions.
\end{proof}

\subsection{Recovery of contextual equation induction}

\begin{theorem}[Equation induction]
For every literal equation $E(y)$, $\mathcal G$ derives
\[
 \frac{\Gamma\vdash E(0)\qquad\Gamma,E(z)\vdash E(Sz)}
      {\Gamma\vdash E(t)},
 \qquad z\text{ fresh for }\Gamma\text{ and the fixed parameters}.
 \tag{EI}\label{eq:guarded-ei}
\]
\end{theorem}
\begin{proof}
Put $J(y)=G_yE(y)$.  We first prove the progress premise
$\Gamma,J(z)\vdash J(Sz)$ without making an invalid contextual use of
NC.  Consider the compound formula
\[
 \Phi(y):=J(y)\Rightarrow_g E(Sy).
\]
Under $\Gamma$, prove $\Phi(0)$ by equation-implication introduction:
ignore its local hypothesis, use the supplied proof of $E(0)$, and
apply the supplied step proof specialized at $0$ to obtain $E(1)$.

For fresh $w$, prove $\Phi(Sw)$ by equation-implication introduction.
Its local hypothesis $J(Sw)$ gives $E(Sw)$ by guard removal, and the
supplied step proof specialized at $Sw$ gives $E(SSw)$.
These two branch proofs have retained context $\Gamma$, independent
of the case variable.  Thus NC legitimately gives $\Phi(z)$ under
$\Gamma$.  Under the local hypothesis $J(z)$, equation-implication
elimination gives $E(Sz)$, hence $J(Sz)$ by the successor guard endpoint.

Atomic progress DI now introduces $\mathsf D_{u,y}J(y)$ with unused
$u$.  DE and the theorem $J(0)$ yield $J(Sw)$, and hence $E(Sw)$, for
a fresh $w$.  Finally NC for $E$ uses the supplied base $E(0)$ and
these successor instances to conclude $E(t)$.

The only induction-like premise used here is the single displayed
guarded progress introduction.  In particular, no predecessor,
coverage disjunction, formula induction, or pre-existing EI is assumed.
\end{proof}

\subsection{Exact translations to and from the inherited calculus}

Let $\mathcal C$ be the inherited Choice-DI calculus: the same scalar
terms, arbitrary formulas $s=t$ and $\mathsf D_{x,y}B$, the same
computation, Use, DE, and transport rules, with case DI at every matrix
and progress DI at literal equations.  Define $T:\mathcal C\to\mathcal G$
on finite formulas by
\[
 T(s=t)=(s=t),\qquad
 T(\mathsf D_{x,y}B)=\mathsf D_{x,y}P(TB).
\tag{Translation}\label{eq:guarded-translation}
\]
Each padding binder uses fresh names.  This is a computable structural
translation, preserves free numerical variables, commutes with
capture-free numerical substitution up to bound-variable renaming,
and translates every paired matrix to a compound matrix.
Translate clauses by
$T\langle A;y.B\rangle=\langle TA;y.TB\rangle$.
Let $J:\mathcal G\to\mathcal C$ be literal inclusion.

\begin{theorem}[Translation of derivations]
\[
 \Gamma\vdash_{\mathcal C}A\Longrightarrow
 T\Gamma\vdash_{\mathcal G}TA,
 \qquad
 \Gamma\vdash_{\mathcal G}A\Longrightarrow
 J\Gamma\vdash_{\mathcal C}JA.
\]
The transformations are effective on finite derivations.
\end{theorem}
\begin{proof}
Inclusion is immediate: every guarded rule is an allowed Choice-DI
rule.  For $T$, induct on a Choice-DI derivation.

Computation is unchanged.  Use commutes with substitution, and
transport uses the translated motive.  A case DI premise translates
to a proof from $TB(t,0)$ to $TB(t,Sz)$.  Under the local assumption
$P(TB)(t,0)$, unpad, apply that translated proof, and pad the result.
Compound case DI then proves the required
$\mathsf D_{x,y}P(TB)$.

For a progress DI with equation matrix $E$, its translated step is
still $\Gamma,E(t,z)\vdash E(t,Sz)$.  To introduce
$\mathsf D_{x,y}L(E(x,y))$, use compound case DI.  Its retained base
$L(E(t,0))$ supplies $E(t,0)$.  The recovered contextual EI, under
that base, gives $E(t,Sz)$, which padding converts to the target
$L(E(t,Sz))$.

For DE, the translated major is $\mathsf D_{x,y}P(TB)$.
It opens the clause
$\langle P(TB)(a,0);y.P(TB)(a,y)\rangle$.
The recursively translated minor expects
$\langle TB(a,0);y.TB(a,y)\rangle$.  Replace every use of this
expected clause by: pad its antecedent, use the opened clause at the
same term, and unpad its consequent.  This effective local proof
replacement preserves the fresh witness condition and the conclusion.
\end{proof}

\begin{lemma}[Compound-matrix congruence]
If canonical compound matrices $B,C$ are uniformly interderivable,
then $\mathsf D_{x,y}B\dashv\vdash\mathsf D_{x,y}C$ in $\mathcal G$.
\end{lemma}
\begin{proof}
Open the source binder at a fresh witness $a$.  Introduce the target
by compound case DI at that witness.  Convert its retained base to
the source base, use the opened source clause at $Sz$, and convert
the successor instance to the target instance.  Close DE.  The
reverse conversion proves the converse.
\end{proof}

\begin{lemma}[The guarded atomic composite]
For an admissible zero-initialized atomic matrix $B(x,y)$,
\[
 \mathsf D_{x,y}B\dashv\vdash\mathsf D_{x,y}L(B)
 \quad\text{in }\mathcal G.
 \tag{Atomic composite}\label{eq:guarded-atomic-composite}
\]
\end{lemma}
\begin{proof}
Forwards, open the source clause at $a$.  Introduce the target by
compound case DI: unpad its base, use the source clause at $Sz$, and
pad the result.

Backwards, open the clause of $\mathsf D_{x,y}L(B)$ at $a$.
The two zero computations prove $B(a,0)$, and padding proves
$L(B(a,0))$.
Thus the opened clause gives $L(B(a,Sz))$, and unpadding gives
$B(a,Sz)$, independently of the local assumption $B(a,z)$.
Atomic progress DI therefore introduces $\mathsf D_{x,y}B$ at
witness $a$.  Close DE.  This argument uses the true base explicitly;
it would fail for an arbitrary unguarded atomic matrix.
\end{proof}

\begin{theorem}[Both composites and bi-interpretability]
For all $\mathcal C$-formulas $A$ and all canonical
$\mathcal G$-formulas $B$,
\[
 A\dashv\vdash_{\mathcal C}JT(A),\qquad
 B\dashv\vdash_{\mathcal G}TJ(B).
\]
Consequently $\mathcal C$ and $\mathcal G$ admit effective structural
intertranslations preserving their full consequence relations, with
both native formula composites and literal identity numerical
translations.
\end{theorem}
\begin{proof}
In $\mathcal C$, padding is interderivable with its argument, and
case introduction proves binder congruence for arbitrary matrices.
Structural induction gives the first displayed equivalence.

For the second, induct on a canonical $B$.  Equations are unchanged.
At a compound matrix, the induction hypothesis and compound-matrix
congruence apply.  At a guarded atomic matrix, the statement is
exactly \eqref{eq:guarded-atomic-composite}.  The equivalences are
uniform under retained assumptions, so they extend to clauses and
sequent contexts by the same local replacement of Use instances.

Both interpretations leave the numerical sort, every term, and
numerical equality untouched.  Both numerical composites therefore
have the identity graph $x=y$; no quotient, retract, numeral restriction,
or extra coding sort is involved.
\end{proof}

For the first-order bi-interpretability statement, use the ordinary
realization that reads every canonical
$\mathsf D_{x,y}B$ uniformly as
$\exists x(B(x,0)\to\forall yB(x,y))$ and leaves numerical equality
and all scalar terms unchanged.  Composing the displayed translations
and composites with the inherited full-domain bridge between
$\mathcal C$ and iPRA yields a bi-interpretation of this ordinary
realization of $\mathcal G$ with full first-order iPRA.
The contradiction $0=1$ is
unchanged, so equiconsistency follows as well.  The inherited arithmetic
bridge supplies the primitive-recursive compiler and the ordinary
logical interpretation; this section does not claim those inherited
constructions as new.

\subsection{Why an unrestricted checker can still be sound}

There is a useful, but strictly separate, administrative calculus
$\mathcal U$.  It permits arbitrary atomic paired matrices and uses
the same root-determined DI, with ordinary Q-style clauses on DE.
It is not claimed that all of $\mathcal U$ is bi-interpretable with iPRA.

Give arbitrary administrative formulas their following ordinary
interpretations:
\[
 (s=t)^\dagger=(s=t),
\]
\[
 (\mathsf D_{x,y}B)^\dagger=
 \begin{cases}
 \exists x\,\forall y\,(B(x,y)\to B(x,Sy)),
       &B\text{ is a literal equation},\\[2pt]
 \exists x\,(B^\dagger(x,0)\to\forall y\,B^\dagger(x,y)),
       &B\text{ is compound}.
 \end{cases}
 \tag{Mixed administrative reading}\label{eq:guarded-mixed}
\]
Retained clauses keep the ordinary reading
$A^\dagger\to\forall y\,B^\dagger$.
If the raw codec also permits administrative $Q$-closures as
formula-shaped objects, including under transport or a paired binder,
extend this interpretation by
$(Q_yB)^\dagger=B^\dagger(0)\to\forall yB^\dagger(y)$.
Such a matrix is compound for root-determined DI.  This adds no
canonical object-formula constructor and preserves the following
soundness argument.

Every $\mathcal U$ inference is sound in iPRA under this reading.
Atomic DI directly proves its progressive reading; compound case DI
uses numerical coverage.  In atomic DE, contextual equation induction
turns $\forall y(E(y)\to E(Sy))$ into
$E(0)\to\forall yE(y)$, exactly the clause that DE supplies.
All other rules are immediate, including transport: numerical
substitution does not change whether a matrix is an equation.

On canonical zero-initialized matrices, the mixed and ordinary readings
agree provably in iPRA.  Indeed $B(x,0)$ is a theorem, so, with $x$ fixed,
\[
 \forall y\,(B(x,y)\to B(x,Sy))
 \quad\longleftrightarrow\quad
 B(x,0)\to\forall y\,B(x,y).
\]
The forward direction is equation induction; the reverse uses the
true base to obtain every instance.  Structural induction extends
this agreement to all canonical guarded formulas.

Thus administrative proofs with unguarded intermediate formulas are
conservative on canonical guarded conclusions: translate an accepted
derivation into iPRA, use the canonical agreement, and translate back
through the established bi-interpretation.  Conversely, every canonical
guarded derivation is already an administrative derivation.

This justifies a checker contract consisting of soundness and
completeness for canonical guarded statements while allowing broader
administrative certificates.  It does \emph{not} make the raw checker
a validator of the guarded public grammar.  It also does not assert
completeness for every noncanonical encoding under the mixed reading.
If a comparison requires rejection of all noncanonical public goals,
the corresponding syntax validator must be implemented and counted.

\subsection{Why merely deleting the choice would not suffice}

Without the guarded canonical language, root-deterministic DI does
not preserve the original native consequence relation.  Let
\[
 p(y)=I(y,0;u.C(u,1,0)),\qquad E(y)\equiv(p(y)=1),
\]
and leave the existential witness coordinate unused.  Computation
gives $p(0)=0,p(1)=1,p(2)=0$.  The original case introduction proves
$\mathsf D_{x,y}E(y)$: its base gives $0=1$, from which equality
transport and selector computation prove every equation.

By contrast, the progressive reading of this formula is false in the
standard natural numbers, because $E(1)$ holds and $E(2)$ does not.
The mixed interpretation proves that it is not a theorem of the
unrestricted root-deterministic calculus.  The guard, the explicit
translations, and both composite proofs are therefore substantive
parts of the reduction, rather than optional notation.
